% The Figure structure intentionally includes the non-floating caption.
% The box must stay unbreakable: inside this minipage a breakable tcolorbox
% still splits against \pagegoal-\pagetotal, and the stacked parts overlap.
\par\noindent\begin{minipage}{\linewidth}
\tagstructbegin{tag=Figure,alttext={三条二维向量按行放置形成三乘二矩阵，转置后按列排列得到二乘三矩阵}}
\begin{tcolorbox}[
  enhanced,breakable=false,colback=paper,colframe=linegray,
  boxrule=.7pt,arc=2mm,left=3mm,right=3mm,top=3mm,bottom=3mm,
  fonttitle=\sffamily\bfseries,coltitle=ink,
  title={图解：同一批向量，可以按行排，也可以按列排}
]
\small
这里先不谈 Attention。假设有三条二维向量：
$(1,0)$、$(0,1)$、$(1,1)$。
\textbf{每种底色始终表示同一条向量}，转置只改变行列方向，不改变数字。

\vspace{1.5mm}
\noindent
\begin{minipage}[t]{.49\linewidth}
\centering
\textbf{按行排：三行、两列}\\[2mm]
\begin{tikzpicture}[x=10mm,y=9mm]
\foreach \r/\shade in {1/paper,2/linegray!25,3/linegray!15}{
  \foreach \c in {1,2}{
    \filldraw[draw=linegray,fill=\shade,line width=.55pt]
      (\c-1,1-\r) rectangle (\c,-\r);
  }
}
\node at (.5,-.5){\sffamily 1}; \node at (1.5,-.5){\sffamily 0};
\node at (.5,-1.5){\sffamily 0}; \node at (1.5,-1.5){\sffamily 1};
\node at (.5,-2.5){\sffamily 1}; \node at (1.5,-2.5){\sffamily 1};
\end{tikzpicture}
\\[1mm]
$A:\ 3\times2$\\[-1mm]
{\footnotesize 每行一条向量}
\end{minipage}%
\begin{minipage}[t]{.49\linewidth}
\centering
\textbf{按列排：两行、三列}\\[2mm]
\begin{tikzpicture}[x=10mm,y=9mm]
\foreach \c/\shade in {1/paper,2/linegray!25,3/linegray!15}{
  \foreach \r in {1,2}{
    \filldraw[draw=linegray,fill=\shade,line width=.55pt]
      (\c-1,1-\r) rectangle (\c,-\r);
  }
}
\node at (.5,-.5){\sffamily 1}; \node at (1.5,-.5){\sffamily 0};
\node at (2.5,-.5){\sffamily 1};
\node at (.5,-1.5){\sffamily 0}; \node at (1.5,-1.5){\sffamily 1};
\node at (2.5,-1.5){\sffamily 1};
\end{tikzpicture}
\\[1mm]
$A^\top:\ 2\times3$\\[-1mm]
{\footnotesize 每列一条向量}
\end{minipage}
\[
A^\top=\begin{bmatrix}1&0&1\\0&1&1\end{bmatrix}.
\]

\textbf{先看方向，再检查形状。}
一般地，$M$ 条 $d$ 维向量，按行排列是
$\underbrace{M\times d}_{\text{一条向量占一行}}$，
转置后按列排列是
$\underbrace{d\times M}_{\text{一条向量占一列}}$。
\end{tcolorbox}

\begin{center}\captionof{figure}{行向量与列向量的转置关系}\label{fig:row-column}\end{center}
\tagstructend
\end{minipage}\par
